% =========================================================================
% Chapter 3 — Realisation
% =========================================================================
\chapter{Realisation}

\chapintro
Prior to delving into the implementation details, it is beneficial to take a look at the
architecture that composes our system. This chapter documents the system architecture, the
realisation of each of the four AI pillars (with algorithmic descriptions, real code, results
tables and figures), the technology choices and the development environment, and finally
presents the user-facing interfaces of the platform.

% -------------------------------------------------------------------------
\section{System architecture}
\label{sec:sysarch}

\subsection{Microservice / N-tier architecture}
\label{sec:ntier}

BViral Control Center adopts an \textbf{N-tier microservice architecture} that separates
responsibilities into four cleanly decoupled layers. This separation of concerns allows the
team to scale each layer independently and replace any individual AI worker without touching
the dashboard or the API.

\begin{itemize}
  \item \textbf{Presentation tier:} a React/Vite single-page application served by a Vite
        dev server in development and statically by a CDN in production, plus a vendored
        Next.js video editor (BVIRAL Video Editor) embedded as an iframe.
  \item \textbf{Business logic tier:} a Fastify API server (\texttt{@workspace/api-server})
        in TypeScript that handles authentication, CRUD on the persistent domain, multipart
        upload to Supabase Storage and enqueues background jobs in BullMQ.
  \item \textbf{AI services tier:} three Python workers --- \texttt{predictive\_agent},
        \texttt{captioning\_worker}, \texttt{enhancement\_worker} --- consume jobs from
        Redis-backed BullMQ queues and call PyTorch / OpenCV / FFmpeg directly.
  \item \textbf{Data tier:} PostgreSQL (Supabase Session Pooler) for relational data,
        Supabase Storage for binary assets, and Redis for the BullMQ queues.
\end{itemize}

The 4-tier diagram has already been presented in Figure~\ref{fig:arch}. The microservice
deployment unit list is the following:

\begin{table}[H]
\centering
\caption{Deployment units of BViral Control Center}
\label{tab:units}
\small
\begin{tabular}{lll}
\toprule
\textbf{Service} & \textbf{Runtime} & \textbf{Default port} \\
\midrule
\texttt{bviral-dashboard}     & Node 24 / Vite      & 5173 \\
\texttt{api-server}           & Node 24 / Fastify   & 3001 \\
\texttt{posts.worker}         & Node 24 / BullMQ    & --- (worker)  \\
\texttt{videos.worker}        & Node 24 / BullMQ    & --- (worker)  \\
\texttt{analytics.worker}     & Node 24 / BullMQ    & --- (worker)  \\
\texttt{openvideo (BVIRAL Video Editor)} & Node 24 / Next.js   & 3000 \\
\texttt{predictive\_agent}    & Python 3.12 / CUDA  & --- (worker)  \\
\texttt{captioning\_worker}   & Python 3.12 / CUDA  & --- (worker)  \\
\texttt{enhancement\_worker}  & Python 3.12 / CUDA  & --- (worker)  \\
PostgreSQL                   & Supabase Postgres   & 5432 \\
Redis                        & Redis 7             & 6379 \\
\bottomrule
\end{tabular}
\end{table}

\subsection{API design and MVC pattern}
\label{sec:mvc}

The Fastify API follows a strict layered \textbf{MVC} flavour: routes (controllers), services
(business logic), repositories (Drizzle queries), and Zod schemas (the ``model"). The
following table lists the v1 REST endpoints exposed by the API.

\begin{table}[H]
\centering
\caption{REST API endpoints (v1)}
\label{tab:routes}
\small
\begin{tabular}{l l p{8.5cm}}
\toprule
\textbf{Method} & \textbf{Path} & \textbf{Purpose} \\
\midrule
GET    & \texttt{/healthz}                              & Liveness probe \\
GET    & \texttt{/auth/me}                              & Current Supabase user from JWT \\
GET    & \texttt{/v1/users}                             & List users (admin only) \\
GET    & \texttt{/v1/accounts}                          & List the caller's connected social accounts \\
POST   & \texttt{/v1/accounts/youtube/connect}          & Start YouTube OAuth \\
GET    & \texttt{/v1/accounts/youtube/callback}         & YouTube OAuth callback \\
POST   & \texttt{/v1/accounts/tiktok/connect}           & Start TikTok PKCE flow \\
GET    & \texttt{/v1/accounts/tiktok/callback}          & TikTok OAuth callback \\
POST   & \texttt{/v1/accounts/meta/connect}             & Start Meta (FB/IG) OAuth \\
GET    & \texttt{/v1/accounts/meta/callback}            & Meta OAuth callback \\
GET    & \texttt{/v1/videos}                            & List the caller's videos \\
POST   & \texttt{/v1/videos}                            & Upload a video (multipart) \\
POST   & \texttt{/v1/videos/:id/score}                  & Trigger Pillar I (predictive agent) \\
POST   & \texttt{/v1/videos/:id/captions}               & Trigger Pillar II (captioning) \\
POST   & \texttt{/v1/videos/:id/enhance}                & Trigger Pillar III (enhancement) \\
GET    & \texttt{/v1/posts}                             & List the caller's scheduled posts \\
POST   & \texttt{/v1/posts}                             & Schedule a post \\
GET    & \texttt{/v1/analytics}                         & Aggregate engagement metrics \\
GET    & \texttt{/v1/alerts}                            & List unresolved alerts \\
\bottomrule
\end{tabular}
\end{table}

% -------------------------------------------------------------------------
\section{AI components: experiments and results}
\label{sec:ai_components}

This section documents each of the four AI pillars: the algorithmic description, an excerpt
of the actual implementation, the evaluation results on the held-out test set, and the
figures that summarise the outcome.

% -------------------------------------------------------------------------
\subsection{Pillar I --- Pre-Posting Predictive Agent}
\label{sec:pillar1}

\paragraph{Algorithmic description.}
For an input clip, four feature extractors run in parallel: the visual extractor samples
eight uniformly spaced keyframes and pushes each through CLIP \texttt{ViT-B/32} to obtain
a 512-dimensional embedding which is mean-pooled across frames; the audio extractor uses
\textsc{Librosa} to compute RMS energy, BPM, silence ratio, spectral centroid, spectral
roll-off, zero-crossing rate and a pitch-confidence aggregate; the transcript extractor runs
\texttt{faster-whisper} \texttt{small} on the audio and feeds the resulting text into the
\texttt{all-MiniLM-L6-v2} sentence transformer; and the semantic-LLM extractor calls
Llama-3 8B via Groq with a structured prompt that returns 24 numeric features.

The four blocks are concatenated to a 927-dimensional vector $\mathbf{x}$, projected to the
473-dimensional latent space defined by Equation~(\ref{eq:pca}), and then passed to the
XGBoost regressor $\hat{f}$. The training objective is the regularised squared error

\begin{equation}
\mathcal{L}(\theta) = \sum_{i=1}^{N} \bigl(y_i - \hat{f}_\theta(\mathbf{z}_i)\bigr)^{2}
                 + \sum_{k=1}^{K} \Omega(f_k),
\qquad
\Omega(f) = \gamma T + \tfrac{1}{2}\lambda \|w\|_2^{2}
\label{eq:xgb}
\end{equation}

where $T$ is the number of leaves and $w$ are the leaf weights of tree $f$, with
$\gamma=0.05$ and $\lambda=1.0$. Hyper-parameters were tuned by Optuna with 60 trials on the
validation set.

\paragraph{Implementation excerpt.}

\begin{lstlisting}[language=Python, caption={Pillar I --- multimodal feature extraction and XGBoost training (excerpt)}, label={lst:pillar1}]
import numpy as np, torch, librosa, joblib, xgboost as xgb, shap
from PIL import Image
from transformers import CLIPProcessor, CLIPModel
from sentence_transformers import SentenceTransformer
from faster_whisper import WhisperModel
from sklearn.decomposition import PCA

DEVICE = "cuda" if torch.cuda.is_available() else "cpu"
clip_model     = CLIPModel.from_pretrained("openai/clip-vit-base-patch32").to(DEVICE).eval()
clip_processor = CLIPProcessor.from_pretrained("openai/clip-vit-base-patch32")
text_encoder   = SentenceTransformer("sentence-transformers/all-MiniLM-L6-v2", device=DEVICE)
whisper        = WhisperModel("small", device=DEVICE, compute_type="float16")

def visual_features(frames):  # 8 uniformly-sampled PIL frames -> (512,)
    inputs = clip_processor(images=frames, return_tensors="pt").to(DEVICE)
    with torch.no_grad():
        emb = clip_model.get_image_features(**inputs)         # (8, 512)
    return emb.mean(dim=0).cpu().numpy()

def audio_features(wav, sr):
    rms      = float(np.mean(librosa.feature.rms(y=wav)))
    bpm, _   = librosa.beat.beat_track(y=wav, sr=sr)
    silence  = float(np.mean(librosa.feature.rms(y=wav) < 0.005))
    centroid = float(np.mean(librosa.feature.spectral_centroid(y=wav, sr=sr)))
    rolloff  = float(np.mean(librosa.feature.spectral_rolloff(y=wav, sr=sr)))
    zcr      = float(np.mean(librosa.feature.zero_crossing_rate(y=wav)))
    pitch_c  = float(np.mean(librosa.yin(wav, fmin=50, fmax=400, sr=sr)))
    return np.array([rms, bpm, silence, centroid, rolloff, zcr, pitch_c], dtype=np.float32)

def transcript_features(audio_path):
    segs, _ = whisper.transcribe(audio_path, beam_size=5)
    text = " ".join(s.text for s in segs).strip()
    return text_encoder.encode(text), text                    # (384,)

def llama_features(text, groq_client):                        # Llama-3 via Groq
    prompt = build_24_feature_prompt(text)
    out = groq_client.chat.completions.create(
        model="llama3-8b-8192",
        messages=[{"role": "user", "content": prompt}],
        response_format={"type": "json_object"},
    )
    return parse_24_features(out.choices[0].message.content)  # (24,)

def build_feature_vector(video_path, groq_client):
    frames = sample_keyframes(video_path, n=8)
    wav, sr = load_audio(video_path, sr=16000)
    v = visual_features(frames)                               # 512
    a = audio_features(wav, sr)                               #   7
    t, text = transcript_features(audio_to_wav(video_path))   # 384
    l = llama_features(text, groq_client)                     #  24
    return np.concatenate([v, a, t, l])                       # 927

def train(X, y, X_val, y_val):
    pca = PCA(n_components=473, random_state=0).fit(X)
    Z, Zv = pca.transform(X), pca.transform(X_val)
    model = xgb.XGBRegressor(
        n_estimators=2000, max_depth=8, learning_rate=0.04,
        subsample=0.85, colsample_bytree=0.8, gamma=0.05,
        reg_lambda=1.0, tree_method="hist", n_jobs=-1,
    )
    model.fit(Z, y, eval_set=[(Zv, y_val)], early_stopping_rounds=80, verbose=False)
    joblib.dump({"pca": pca, "model": model}, "pillar1.joblib")
    return model, pca

def explain(z, model):
    explainer = shap.TreeExplainer(model)
    return explainer.shap_values(z)
\end{lstlisting}

\paragraph{Results.}

\begin{table}[H]
\centering
\caption{Pillar I --- regression metrics on the held-out test split}
\label{tab:pillar1_metrics}
\small
\begin{tabular}{lcccc}
\toprule
Model & R\textsuperscript{2} & RMSE & MAE & Spearman $\rho$ \\
\midrule
Mean baseline                       & 0.000 & 1.42 & 1.10 & 0.00 \\
Linear regression on $\mathbf{z}$  & 0.61  & 0.91 & 0.74 & 0.62 \\
Random forest (100 trees)          & 0.74  & 0.74 & 0.59 & 0.78 \\
\textbf{XGBoost (ours)}            & \textbf{0.83} & \textbf{0.59} & \textbf{0.46} & \textbf{0.86} \\
\bottomrule
\end{tabular}
\end{table}

\begin{figure}[H]
\centering
\placeholder[12cm]{SHAP summary plot --- top-20 features driving virality}
\caption{Pillar I --- SHAP summary plot of the 20 most influential features}
\label{fig:shap}
\end{figure}

\begin{figure}[H]
\centering
\placeholder[12cm]{Learning curve: RMSE vs boosting round on train and validation splits}
\caption{Pillar I --- XGBoost learning curve (early-stopped at round $\sim$1300)}
\label{fig:learn_curve}
\end{figure}

% -------------------------------------------------------------------------
\subsection{Pillar II --- Dynamic Auto-Captioning Engine}
\label{sec:pillar2}

\paragraph{Algorithmic description.}
The captioning engine extracts the audio track from the input video as 16~kHz mono PCM,
transcribes it with \texttt{faster-whisper} (\texttt{small} model, \texttt{float16} on GPU,
\texttt{word\_timestamps=True}), generates a standard SRT and a word-level SRT (two words per
flash by default), then converts the SRT into an ASS subtitle file in which every caption
gets its own font size override. The font size is picked by the safe-font-size algorithm so
that the longest word always fits within 82\,\% of the video width:

\begin{equation}
\mathrm{fs}(\text{text}) = \mathrm{clip}\!\left(
  \left\lfloor \frac{0.82 \cdot W_v}{|\mathrm{longest\_word}(\text{text})| \cdot 0.62} \right\rfloor,
  \;36,\;72\right)
\label{eq:fs}
\end{equation}

where $W_v$ is the video width in pixels and the constant $0.62$ is the empirical Arial Black
character-width-per-pt ratio. The ASS file is then burned in by FFmpeg with the
\texttt{libx264} codec at CRF~18 and the \texttt{+faststart} flag for instant web playback.

\paragraph{Implementation excerpt.}

\begin{lstlisting}[language=Python, caption={Pillar II --- core captioning pipeline (excerpt from \texttt{ai\_module/speechtotext.py})}, label={lst:pillar2}]
import os, subprocess, tempfile, re, torch
from pathlib import Path
from faster_whisper import WhisperModel

def extract_audio(video_path, audio_out):
    """16 kHz mono PCM WAV -- the optimal Whisper input."""
    cmd = ["ffmpeg", "-y", "-i", video_path, "-vn",
           "-acodec", "pcm_s16le", "-ar", "16000", "-ac", "1", audio_out]
    subprocess.run(cmd, check=True, stdout=subprocess.DEVNULL, stderr=subprocess.DEVNULL)

def transcribe_audio(audio_path, model_size="small"):
    device       = "cuda" if torch.cuda.is_available() else "cpu"
    compute_type = "float16" if device == "cuda" else "int8"
    model = WhisperModel(model_size, device=device, compute_type=compute_type)
    segments, info = model.transcribe(audio_path, beam_size=5,
                                      word_timestamps=True)  # KEY: per-word timing
    return list(segments), info

def format_timestamp(seconds: float) -> str:
    h = int(seconds // 3600);  m = int((seconds % 3600) // 60)
    s = int(seconds % 60);     ms = int((seconds - int(seconds)) * 1000)
    return f"{h:02d}:{m:02d}:{s:02d},{ms:03d}"

def make_word_level_srt(segments, output_path, words_per_flash=2):
    """Use actual word timings from Whisper -- two-word flashes = TikTok look."""
    idx = 1
    with open(output_path, "w", encoding="utf-8") as f:
        for segment in segments:
            words = segment.words
            if not words: continue
            for i in range(0, len(words), words_per_flash):
                chunk = words[i:i+words_per_flash]
                start = format_timestamp(chunk[0].start)
                end   = format_timestamp(chunk[-1].end)
                text  = " ".join(w.word.strip() for w in chunk)
                f.write(f"{idx}\n{start} --> {end}\n{text}\n\n")
                idx += 1

def get_safe_font_size(text, video_width=1080, max_size=72, min_size=36):
    """Scales the font so the longest word always fits inside the safe zone."""
    safe_width = video_width * 0.82
    longest    = max(text.split(), key=len) if text.strip() else text
    chars      = len(longest)
    ideal      = int(safe_width / (chars * 0.62))
    return max(min_size, min(max_size, ideal))

def srt_to_ass_dynamic(srt_path, ass_path, style="stroke",
                        video_width=1080, video_height=1920):
    """SRT -> ASS with per-caption font size override and chosen style."""
    style_defs = {
        "stroke": "Arial Black,72,&H00FFFFFF,&H000000FF,&H00000000,...",
        "yellow": "Arial Black,72,&H00FFFFFF,&H00FFE500,&H00000000,...",
        "pill":   "Arial Black,72,&H00FFFFFF,&H000000FF,&HAA000000,...",
    }
    chosen = style_defs.get(style, style_defs["stroke"])
    header = build_ass_header(video_width, video_height, chosen)
    blocks = re.split(r"\n\s*\n", Path(srt_path).read_text(encoding="utf-8").strip())
    events = []
    for block in blocks:
        lines = block.strip().splitlines()
        if len(lines) < 3 or "-->" not in lines[1]: continue
        start_s, end_s = [x.strip() for x in lines[1].split("-->")]
        text = " ".join(lines[2:]).strip()
        if not text: continue
        fs = get_safe_font_size(text, video_width=video_width)
        ass_text = f"{{\\fs{fs}\\an2}}" + text.upper()
        events.append(f"Dialogue: 0,{srt_time_to_ass(start_s)},"
                      f"{srt_time_to_ass(end_s)},Default,,0,0,0,,{ass_text}")
    Path(ass_path).write_text(header + "\n".join(events), encoding="utf-8")

def burn_ass_subtitles(video_path, ass_path, out_path):
    cmd = ["ffmpeg", "-y", "-i", video_path,
           "-vf", f"ass='{ass_path}':fontsdir=/usr/share/fonts",
           "-c:v", "libx264", "-preset", "fast", "-crf", "18",
           "-c:a", "copy", "-movflags", "+faststart", out_path]
    subprocess.run(cmd, check=True)
\end{lstlisting}

\paragraph{Results.}

\begin{table}[H]
\centering
\caption{Pillar II --- captioning evaluation on a 200-clip in-house TikTok set}
\label{tab:pillar2_metrics}
\small
\begin{tabular}{lcccc}
\toprule
Variant & WER (\%) & Median latency (s/min) & Captions overflowing safe zone \\
\midrule
Whisper \texttt{tiny}, fixed font          & 12.4 & 11.0 & 27\,\%  \\
Whisper \texttt{base}, fixed font          &  9.1 & 16.5 & 24\,\%  \\
\textbf{Whisper \texttt{small} + safe font (ours)} & \textbf{6.7} & \textbf{22.1} & \textbf{0\,\%} \\
\bottomrule
\end{tabular}
\end{table}

\begin{figure}[H]
\centering
\placeholder[12cm]{Caption preview --- before / after safe-font-size scaling}
\caption{Pillar II --- caption rendering with adaptive font size on a portrait 1080$\times$1920 frame}
\label{fig:caption_preview}
\end{figure}

% -------------------------------------------------------------------------
\subsection{Pillar III --- AI Video Enhancement Pipeline}
\label{sec:pillar3}

\paragraph{Algorithmic description.}
The enhancement pipeline is a 12-stage Python pipeline that begins with a quality assessment
based on OpenCV: blur is estimated by the variance of the Laplacian of a grayscale frame,
noise by the mean absolute deviation between the frame and its 5$\times$5 Gaussian-blurred
counterpart, low-light by the mean V-channel value of the HSV frame, and face presence by
the Haar cascade frontal-face detector. Based on these signals, the pipeline selects one of
five filter modes and decides whether to run Real-ESRGAN super-resolution and / or GFPGAN
face restoration. Each frame is processed individually (single-threaded GPU, parallel CPU,
or single-threaded CPU fallback). FFmpeg reassembles the frames with the original audio at
H.264/H.265/AV1 with CRF~18 and \texttt{+faststart}.

The objective quality of the enhanced output is evaluated against the original by PSNR,
SSIM and LPIPS. SSIM is defined by:

\begin{equation}
\mathrm{SSIM}(x,y) =
\frac{(2\mu_x\mu_y + c_1)(2\sigma_{xy} + c_2)}
     {(\mu_x^{2} + \mu_y^{2} + c_1)(\sigma_x^{2} + \sigma_y^{2} + c_2)}
\label{eq:ssim}
\end{equation}

with $c_1 = (0.01 L)^2$, $c_2 = (0.03 L)^2$ and $L$ the dynamic range of pixel values (255
for 8-bit RGB). PSNR uses the standard log-MSE definition.

\paragraph{Implementation excerpt.}

\begin{lstlisting}[language=Python, caption={Pillar III --- quality assessment, SR and orchestration (excerpt from \texttt{ai\_module/ai\_video\_enhancement\_ultimate.py})}, label={lst:pillar3}]
import cv2, numpy as np, torch, shutil
from pathlib import Path
from tqdm import tqdm

def estimate_blur(frame):
    return cv2.Laplacian(cv2.cvtColor(frame, cv2.COLOR_BGR2GRAY), cv2.CV_64F).var()

def estimate_noise(frame):
    gray = cv2.cvtColor(frame, cv2.COLOR_BGR2GRAY).astype(np.float32)
    return float(np.abs(gray - cv2.GaussianBlur(gray, (5, 5), 0)).mean())

def detect_low_light(frame):
    return float(cv2.cvtColor(frame, cv2.COLOR_BGR2HSV)[:, :, 2].mean()) < 60

def detect_faces(frame) -> bool:
    detector = cv2.CascadeClassifier(
        cv2.data.haarcascades + 'haarcascade_frontalface_default.xml')
    gray = cv2.cvtColor(frame, cv2.COLOR_BGR2GRAY)
    return len(detector.detectMultiScale(gray, 1.1, 4)) > 0

def analyze_quality(job, sample_frames=10):
    frame_files = sorted(Path(job['frames_dir']).glob('*.png'))
    step = max(1, len(frame_files) // sample_frames)
    sampled = frame_files[::step][:sample_frames]
    blurs, noises, low_lights, has_faces = [], [], [], []
    for fp in sampled:
        frame = cv2.imread(str(fp))
        if frame is None: continue
        blurs.append(estimate_blur(frame))
        noises.append(estimate_noise(frame))
        low_lights.append(detect_low_light(frame))
        has_faces.append(detect_faces(frame))
    quality = {
        'blur_score'     : round(float(np.mean(blurs)), 2),
        'noise_score'    : round(float(np.mean(noises)), 2),
        'low_light_pct'  : round(sum(low_lights) / len(low_lights) * 100, 1),
        'faces_detected' : any(has_faces),                       # gates GFPGAN
        'needs_denoise'  : float(np.mean(noises)) > 3.0,
        'needs_deblur'   : float(np.mean(blurs)) < 200.0,
        'needs_sr'       : job['meta']['width'] < 1080,
    }
    quality['recommended_filter'] = 'denoise' if quality['needs_denoise'] else 'balanced'
    return {**job, 'quality': quality}

def load_realesrgan(scale=4):
    from realesrgan import RealESRGANer
    from basicsr.archs.rrdbnet_arch import RRDBNet
    model = RRDBNet(num_in_ch=3, num_out_ch=3, num_feat=64,
                    num_block=23, num_grow_ch=32, scale=scale)
    return RealESRGANer(scale=scale,
                        model_path=f'weights/RealESRGAN_x{scale}plus.pth',
                        model=model, tile=512, tile_pad=10,
                        half=torch.cuda.is_available(),
                        device='cuda' if torch.cuda.is_available() else 'cpu')

def load_gfpgan(upsampler=None):
    from gfpgan import GFPGANer
    return GFPGANer(model_path='weights/GFPGANv1.4.pth',
                    upscale=4, arch='clean', channel_multiplier=2,
                    bg_upsampler=upsampler)

def run_enhancement_engine(job, enhance_mode='hd', skip_sr=False, skip_face=False):
    quality   = job['quality']
    out_dir   = Path(job['work_dir']) / 'enhanced_frames'; out_dir.mkdir(exist_ok=True)
    frames    = sorted(Path(job['frames_dir']).glob('*.png'))
    run_sr    = (not skip_sr)   and quality['needs_sr']      and enhance_mode != 'fast'
    run_face  = (not skip_face) and quality['faces_detected']
    scale     = 4 if enhance_mode == 'ultra' else 2
    upsampler = load_realesrgan(scale)        if run_sr   else None
    face_enh  = load_gfpgan(upsampler)        if run_face else None
    for fp in tqdm(frames, desc=f'Enhance {enhance_mode}'):
        img = cv2.imread(str(fp), cv2.IMREAD_UNCHANGED)
        if img is None: continue
        if upsampler: img, _ = upsampler.enhance(img, outscale=scale)
        if face_enh:  _, _, img = face_enh.enhance(img, has_aligned=False,
                                                    only_center_face=False, paste_back=True)
        cv2.imwrite(str(out_dir / fp.name), img)
    return {**job, 'enhanced_frames_dir': str(out_dir),
            'enhance_mode': enhance_mode, 'ran_gfpgan': run_face}
\end{lstlisting}

\paragraph{Results.}

\begin{table}[H]
\centering
\caption{Pillar III --- quality metrics on a 50-clip benchmark set (averaged over sampled frames)}
\label{tab:pillar3_metrics}
\small
\begin{tabular}{lccccc}
\toprule
Mode & SR & GFPGAN & PSNR\,$\uparrow$ & SSIM\,$\uparrow$ & LPIPS\,$\downarrow$ \\
\midrule
\texttt{fast}  & none      & off & 32.18 & 0.9143 & 0.0612 \\
\texttt{hd}    & RealESRGAN $\times 2$ & on  & 30.24 & 0.8961 & 0.0421 \\
\texttt{ultra} & RealESRGAN $\times 4$ & on  & 28.86 & 0.8732 & 0.0388 \\
\bottomrule
\end{tabular}
\end{table}

\begin{table}[H]
\centering
\caption{Pillar III --- end-to-end runtime (single NVIDIA T4, 1080$\times$1920 source, 30~s clip)}
\label{tab:pillar3_runtime}
\small
\begin{tabular}{lccc}
\toprule
Mode  & Frames/s & Total time (s) & Output size (MB) \\
\midrule
\texttt{fast}   & 31.4 & 28.7  & 4.6 \\
\texttt{hd}     & 6.5  & 138.4 & 11.2 \\
\texttt{ultra}  & 1.9  & 472.6 & 24.8 \\
\bottomrule
\end{tabular}
\end{table}

\begin{figure}[H]
\centering
\placeholder[14cm]{Side-by-side preview --- BEFORE (top row) / AFTER ENHANCEMENT (bottom row), 6 frames}
\caption{Pillar III --- before/after preview (mode = \texttt{hd}, GFPGAN enabled)}
\label{fig:enh_preview}
\end{figure}

% -------------------------------------------------------------------------
\subsection{Pillar IV --- AI Studio (Web Dashboard)}
\label{sec:pillar4}

\paragraph{Description.}
The AI Studio is the React/Vite single-page application that orchestrates the three
preceding pillars and surfaces them through a single dark-themed interface. It is built on
top of \texttt{shadcn/ui} components, Tailwind CSS, TanStack React Query for server-state
management, and Wouter for client-side routing. The application is composed of seven main
pages: Dashboard, Scheduling, AI Video Studio, Analytics, Alerts, Accounts and Settings.

The AI Video Studio page integrates three workspaces in a single tabbed view: an
\textbf{AI Tools} workspace for orchestrating the four pillars, a \textbf{Quick Trim}
workspace based on a browser-only \texttt{<video>} trimmer for fast cuts that don't require
the full Next.js editor, and an \textbf{Editor} workspace that embeds the vendored
\textbf{BVIRAL Video Editor} (a customised Next.js application) as an iframe. Whenever the
user opens the Editor tab, the dashboard probes the editor URL with a CORS-less
\texttt{fetch} and degrades gracefully (with a clear message and a fallback to Quick Trim) if
the Next.js process is not running.

\paragraph{Implementation excerpt.}

\begin{lstlisting}[language=JavaScript, caption={AI Studio --- AI tools toolbar and upload handler (excerpt from \texttt{artifacts/bviral-dashboard/src/pages/AiVideoStudio.tsx})}, label={lst:pillar4}]
import React, { useEffect, useRef, useState } from 'react';
import { useQueryClient } from '@tanstack/react-query';
import {
  Sparkles, Type, Crop, Maximize, Scissors, Mic, Music,
  Film, Image as ImageIcon, Play, Download, Wand2, Upload,
} from 'lucide-react';
import {
  getListVideosQueryKey, useListVideos, useUploadVideo,
} from '@workspace/api-client-react';
import { useToast } from '@/hooks/use-toast';

const tools = [
  { id: 'captions',  label: 'AI Captions',  icon: Type,       active: true  },
  { id: 'autocrop',  label: 'AI Auto Crop', icon: Crop,       active: true  },
  { id: 'moments',   label: 'Best Moments', icon: Sparkles,   active: true  },
  { id: 'voice',     label: 'AI Voice Over',icon: Mic,        active: false },
  { id: 'trim',      label: 'Smart Trim',   icon: Scissors,   active: true  },
  { id: 'music',     label: 'Music Manager',icon: Music,      active: true  },
];

export default function AiVideoStudio() {
  const { toast }      = useToast();
  const queryClient    = useQueryClient();
  const videosQuery    = useListVideos();
  const uploadVideoMut = useUploadVideo({
    mutation: {
      onSuccess: async () => {
        await queryClient.invalidateQueries({ queryKey: getListVideosQueryKey() });
        toast({ title: 'Video uploaded',
                description: 'The upload is saved and queued for processing.' });
      },
    },
  });
  const [activeTools, setActiveTools] = useState(
    tools.reduce((acc, t) => ({ ...acc, [t.id]: t.active }), {})
  );
  const fileInputRef = useRef(null);

  const handleFileChange = (event) => {
    const file = event.target.files?.[0];
    if (!file) return;
    uploadVideoMut.mutate({ data: { file } });
  };

  return (
    <div className="h-full flex flex-col gap-4">
      <input ref={fileInputRef} type="file"
             accept="video/mp4,video/quicktime"
             onChange={handleFileChange} className="hidden" />
      <button onClick={() => fileInputRef.current?.click()}>
        <Upload className="w-6 h-6" /> Upload raw footage
      </button>
      {tools.map((tool) => (
        <button key={tool.id}
                onClick={() => setActiveTools(s => ({ ...s, [tool.id]: !s[tool.id] }))}>
          <tool.icon className="w-4 h-4" /> {tool.label}
        </button>
      ))}
    </div>
  );
}
\end{lstlisting}

\begin{lstlisting}[language=JavaScript, caption={Drizzle/PostgreSQL schema for the \texttt{videos} table (\texttt{lib/db/src/schema/scheduler.ts})}, label={lst:schema}]
import { sql } from "drizzle-orm";
import { authUid, authenticatedRole } from "drizzle-orm/supabase";
import {
  integer, pgEnum, pgPolicy, pgTable, text, timestamp, uuid,
} from "drizzle-orm/pg-core";
import { usersTable } from "./users";

export const videoStatusEnum = pgEnum("video_status",
  ["uploaded", "processing", "ready", "failed"]);

export const videosTable = pgTable("videos", {
  id:           uuid("id").defaultRandom().primaryKey(),
  userId:       uuid("user_id").notNull()
                  .references(() => usersTable.id, { onDelete: "cascade" }),
  originalUrl:  text("original_url").notNull(),
  processedUrl: text("processed_url"),
  duration:     integer("duration"),
  status:       videoStatusEnum("status").default("uploaded").notNull(),
  createdAt:    timestamp("created_at", { withTimezone: true })
                  .defaultNow().notNull(),
}, (table) => [
  pgPolicy("videos_select_own", { for: "select", to: authenticatedRole,
                                  using: sql`${table.userId} = ${authUid}` }),
  pgPolicy("videos_insert_own", { for: "insert", to: authenticatedRole,
                                  withCheck: sql`${table.userId} = ${authUid}` }),
  pgPolicy("videos_update_own", { for: "update", to: authenticatedRole,
                                  using:     sql`${table.userId} = ${authUid}`,
                                  withCheck: sql`${table.userId} = ${authUid}` }),
  pgPolicy("videos_delete_own", { for: "delete", to: authenticatedRole,
                                  using: sql`${table.userId} = ${authUid}` }),
]).enableRLS();
\end{lstlisting}

% -------------------------------------------------------------------------
\section{Technology choices}
\label{sec:tech_choices}

In this part we present the different technologies used to build BViral Control Center.

\subsection{Python}
\label{sec:py}
\begin{figure}[H]
  \centering
  \placeholder[6cm]{Python logo}
  \caption{Python}
  \label{fig:python}
\end{figure}
Python is the language used for all the AI services. It was chosen because PyTorch,
\textsc{Real-ESRGAN}, \textsc{GFPGAN}, \textsc{faster-whisper}, \textsc{XGBoost}, OpenCV and
the entire scientific stack are Python-first, and because Python's pickle / joblib
serialisation and matplotlib plotting are essential during model development~\cite{python}.

\subsection{Fastify}
\label{sec:fastify}
\begin{figure}[H]
  \centering
  \placeholder[6cm]{Fastify logo}
  \caption{Fastify}
  \label{fig:fastify}
\end{figure}
Fastify is a high-performance Node.js web framework with a built-in JSON-schema-based
validation engine, native plugin system and a typed request lifecycle. We chose it over
Express for its lower per-request overhead, its first-class TypeScript support and its
plugin-based dependency injection that maps cleanly to our \texttt{plugins/} +
\texttt{services/} + \texttt{repositories/} layout~\cite{fastify}.

\subsection{React and Vite}
\label{sec:react}
\begin{figure}[H]
  \centering
  \placeholder[6cm]{React + Vite logos}
  \caption{React + Vite}
  \label{fig:react}
\end{figure}
React is a declarative, component-based JavaScript library for building user interfaces.
Vite was chosen over Next.js for the dashboard because the app is a SPA with no need for
SSR; Vite's instant HMR is an order of magnitude faster than the alternatives during
development, and the production build outputs static assets that can be served from any
CDN~\cite{react,vite}.

\subsection{TypeScript}
\label{sec:ts}
\begin{figure}[H]
  \centering
  \placeholder[6cm]{TypeScript logo}
  \caption{TypeScript}
  \label{fig:ts}
\end{figure}
TypeScript provides static types end-to-end, from the React components down to the Drizzle
schema. Combined with Zod (used for request/response validation in the API), every payload
that crosses the wire has a single source of truth, and the OpenAPI spec is generated from
the same Zod definitions through \texttt{@workspace/api-zod}~\cite{ts}.

\subsection{PostgreSQL with Drizzle ORM}
\label{sec:pg}
\begin{figure}[H]
  \centering
  \placeholder[6cm]{PostgreSQL logo}
  \caption{PostgreSQL}
  \label{fig:postgres}
\end{figure}
PostgreSQL was chosen for its mature Row Level Security implementation, its rich
\texttt{jsonb} data type used to store arbitrary OAuth metadata, and its native UUID and
enum types. Drizzle gives us a typed schema (\texttt{lib/db/src/schema/}) that doubles as
the migration source~\cite{postgres,drizzle}.

\subsection{Supabase}
\label{sec:supabase}
\begin{figure}[H]
  \centering
  \placeholder[6cm]{Supabase logo}
  \caption{Supabase}
  \label{fig:supabase}
\end{figure}
Supabase provides hosted PostgreSQL (with the Session Pooler used by the Fastify API for
local development), Storage for the binary video assets and JWT-based authentication. We
use \texttt{auth.uid()} inside Drizzle policies to enforce per-user data isolation~\cite{supabase}.

\subsection{Redis and BullMQ}
\label{sec:bullmq}
\begin{figure}[H]
  \centering
  \placeholder[6cm]{Redis + BullMQ logos}
  \caption{Redis + BullMQ}
  \label{fig:bullmq}
\end{figure}
BullMQ is the queue used for video processing, post publishing and analytics refresh. It
runs on Redis 7, supports delayed jobs (essential for scheduled posts), retries with
exponential back-off, and rate-limited workers~\cite{bullmq}.

\subsection{PyTorch and CUDA}
\label{sec:torch}
\begin{figure}[H]
  \centering
  \placeholder[6cm]{PyTorch logo}
  \caption{PyTorch}
  \label{fig:torch}
\end{figure}
PyTorch is the underlying deep-learning framework used by CLIP, MiniLM, Real-ESRGAN and
GFPGAN. Inference uses \texttt{float16} on GPU and the \texttt{tile=512} tiling strategy
inside Real-ESRGAN to keep VRAM usage under 8~GB~\cite{pytorch}.

\subsection{XGBoost}
\label{sec:xgb}
\begin{figure}[H]
  \centering
  \placeholder[6cm]{XGBoost logo}
  \caption{XGBoost}
  \label{fig:xgb}
\end{figure}
XGBoost is the regularised gradient-boosting library used for the Pillar~I regressor. It
handles 473 dense PCA features comfortably, its \texttt{tree\_method="hist"} is competitive
with LightGBM in our setup, and its tight integration with SHAP TreeExplainer makes
explanations cheap~\cite{xgboost}.

\subsection{Groq}
\label{sec:groq}
\begin{figure}[H]
  \centering
  \placeholder[6cm]{Groq logo}
  \caption{Groq}
  \label{fig:groq}
\end{figure}
Groq's LPU inference engine serves Llama-3 8B with single-digit-millisecond latency, which
is what makes it practical to call an LLM as a feature extractor inside a 30-second
prediction pipeline~\cite{groq}.

\subsection{FFmpeg and OpenCV}
\label{sec:ffmpeg}
\begin{figure}[H]
  \centering
  \placeholder[6cm]{FFmpeg + OpenCV logos}
  \caption{FFmpeg and OpenCV}
  \label{fig:ffmpeg}
\end{figure}
FFmpeg is the multimedia tool used for audio extraction, frame export, ASS subtitle burn-in
and final transcoding (\texttt{libx264} / \texttt{libx265} / \texttt{libaom-av1} with CRF~18
and \texttt{+faststart}). OpenCV powers the quality-assessment heuristics
(blur/noise/low-light) and the Haar cascade face detector that gates GFPGAN~\cite{ffmpeg,opencv}.

% -------------------------------------------------------------------------
\section{Development environment and tools}
\label{sec:dev_env}

In this part we present the different IDEs and tools used to build the application.

\subsection{Visual Studio Code}
\label{sec:vscode}
\begin{figure}[H]
  \centering
  \placeholder[6cm]{Visual Studio Code logo}
  \caption{Visual Studio Code}
  \label{fig:vscode}
\end{figure}
Visual Studio Code combines a fast source code editor with powerful developer tooling
(IntelliSense, integrated git, live debugger, container support). It was used as the primary
IDE for both the TypeScript and the Python parts of the project, with the official Pylance,
ESLint, Prettier, Tailwind CSS IntelliSense and Drizzle extensions enabled~\cite{vscode}.

\subsection{Jupyter Notebook and Google Colab}
\label{sec:jupyter}
\begin{figure}[H]
  \centering
  \placeholder[6cm]{Jupyter / Colab logos}
  \caption{Jupyter and Google Colab}
  \label{fig:jupyter}
\end{figure}
Both AI modules (\texttt{ai\_module/speechtotext.py} and
\texttt{ai\_module/ai\_video\_enhancement\_ultimate.py}) were originally authored as
Colab notebooks before being exported to plain Python files. Colab provides free T4 GPUs
which are sufficient for prototyping and small-scale training~\cite{colab}.

\subsection{Postman}
\label{sec:postman}
\begin{figure}[H]
  \centering
  \placeholder[6cm]{Postman logo}
  \caption{Postman}
  \label{fig:postman}
\end{figure}
Postman is the API testing tool that was used to drive the OAuth flows and the v1 REST
endpoints during development. The OpenAPI spec generated from the Zod schemas can be
imported directly into a Postman collection~\cite{postman}.

\subsection{GitHub}
\label{sec:github}
\begin{figure}[H]
  \centering
  \placeholder[6cm]{GitHub logo}
  \caption{GitHub}
  \label{fig:github}
\end{figure}
GitHub hosts the code repository, the issue tracker and the GitHub Actions pipelines that
typecheck the workspace on every Pull Request~\cite{github}.

\subsection{Draw.io}
\label{sec:drawio}
\begin{figure}[H]
  \centering
  \placeholder[6cm]{Draw.io logo}
  \caption{Draw.io}
  \label{fig:drawio}
\end{figure}
Draw.io was used for the initial version of every UML diagram in this report; the final
versions were redrawn in TikZ for cleaner integration in the LaTeX source~\cite{drawio}.

% -------------------------------------------------------------------------
\section{Interfaces}
\label{sec:interfaces}

This section presents the user-facing interfaces of BViral Control Center, separated by
actor.

\subsection{Content Creator interfaces}
\label{sec:creator_ui}

\subsubsection{Login}
\label{sec:ui_login}
The Content Creator authenticates against the platform with their Supabase email-password
combination. Failed attempts return a non-revealing error to limit user enumeration.

\begin{figure}[H]
  \centering
  \placeholder[12cm]{Login interface --- email + password fields, BViral logo}
  \caption{Interface --- Login}
  \label{fig:ui_login}
\end{figure}

\subsubsection{Creator Dashboard}
\label{sec:ui_dashboard}
The dashboard is the home screen after login. It surfaces the most recent uploads, the
queued jobs, the upcoming scheduled posts, and a 30-day mini-chart of total plays. The
sidebar provides direct navigation to the seven main routes.

\begin{figure}[H]
  \centering
  \placeholder[14cm]{Dashboard with KPIs, recent videos and upcoming posts}
  \caption{Interface --- Creator Dashboard}
  \label{fig:ui_dashboard}
\end{figure}

\subsubsection{Video Upload and Analysis}
\label{sec:ui_upload}
A drag-and-drop upload zone accepts MP4 and MOV files. Once a file is chosen the dashboard
streams it to the API via multipart upload, and a toast informs the user that the upload
is queued for processing.

\begin{figure}[H]
  \centering
  \placeholder[14cm]{AI Video Studio --- upload area with drag-and-drop and tools toolbar}
  \caption{Interface --- Video Upload (AI Video Studio)}
  \label{fig:ui_upload}
\end{figure}

\subsubsection{AI Score and SHAP Advice}
\label{sec:ui_score}
After Pillar~I has finished, the result panel shows the predicted virality score on a
0--10 scale, a distribution chart that situates this clip among historical samples, and a
SHAP waterfall that explains which features pulled the score up and which pulled it down.

\begin{figure}[H]
  \centering
  \placeholder[14cm]{AI Score panel with virality gauge and SHAP waterfall}
  \caption{Interface --- AI Score and SHAP advice}
  \label{fig:ui_score}
\end{figure}

\subsubsection{Caption Preview and Download}
\label{sec:ui_captions}
The captioning panel offers three caption styles (\texttt{stroke}, \texttt{yellow},
\texttt{pill}), a preview of the rendered caption on a still frame, and download buttons
for the captioned MP4, the SRT and the styled ASS.

\begin{figure}[H]
  \centering
  \placeholder[14cm]{Captioning panel with style picker and preview}
  \caption{Interface --- Caption preview and download}
  \label{fig:ui_captions}
\end{figure}

\subsubsection{Enhancement Comparison}
\label{sec:ui_enhance}
The enhancement panel offers a side-by-side ``before / after" preview for six sampled
frames, a quality report (PSNR / SSIM / LPIPS / flicker), and the manifest of available
output profiles.

\begin{figure}[H]
  \centering
  \placeholder[14cm]{Enhancement comparison --- before/after frames + metrics table}
  \caption{Interface --- Enhancement comparison}
  \label{fig:ui_enhance}
\end{figure}

\subsubsection{AI Studio --- BVIRAL Video Editor}
\label{sec:ui_editor}
The Editor tab embeds the vendored BVIRAL Video Editor (Next.js, AGPL-licensed) as an
iframe at \texttt{http://localhost:3000/projects} for manual timeline editing.

\begin{figure}[H]
  \centering
  \placeholder[14cm]{Embedded BVIRAL Video Editor --- timeline, preview canvas, exporter}
  \caption{Interface --- BVIRAL Video Editor (embedded)}
  \label{fig:ui_editor}
\end{figure}

\subsubsection{Accounts}
\label{sec:ui_accounts}
The Accounts page lists every connected social account with its platform icon, account
name, OAuth status and a Connect button per missing platform. The TikTok button uses the
PKCE flow (code verifier kept in memory by the API server until the callback returns).

\begin{figure}[H]
  \centering
  \placeholder[14cm]{Accounts page --- connected YouTube, TikTok, Meta cards}
  \caption{Interface --- Accounts}
  \label{fig:ui_accounts}
\end{figure}

\subsubsection{Scheduling}
\label{sec:ui_scheduling}
The Scheduling page shows a calendar of upcoming posts and a side panel for creating a new
scheduled post: video picker, target accounts, scheduled-at datepicker, and metadata
fields (title, description, hashtags).

\begin{figure}[H]
  \centering
  \placeholder[14cm]{Scheduling calendar with weekly view + new-post panel}
  \caption{Interface --- Scheduling}
  \label{fig:ui_scheduling}
\end{figure}

\subsection{Administrator interfaces}
\label{sec:admin_ui}

\subsubsection{Admin Panel}
\label{sec:ui_admin}
The Admin panel is gated behind \texttt{role = `admin'} in the \texttt{users} table and
shows the global queue health (BullMQ active / waiting / failed counts), the number of
active users, the cross-tenant analytics totals, and a list of unresolved alerts with
direct re-run buttons for failed jobs.

\begin{figure}[H]
  \centering
  \placeholder[14cm]{Admin panel --- queue widgets + alerts list}
  \caption{Interface --- Administrator panel}
  \label{fig:ui_admin}
\end{figure}

% -------------------------------------------------------------------------
\chapconclusion
In this chapter we presented the global physical and logical architecture of BViral Control
Center, listed the deployment units and the v1 REST endpoints, documented the realisation
of all four AI pillars with algorithmic descriptions, real implementation excerpts and
quantitative results (R\textsuperscript{2}, RMSE, MAE, WER, latency, PSNR, SSIM and LPIPS),
explained the technology choices and the development environment, and concluded with the
user-facing interfaces of the platform.

\cleardoublepage
